\subsection{单位根的本原根}

设 \(n \geqslant 1\) 为一个整数。所谓 \(n\) 的欧拉指示函数（记作 \(\varphi(n)\) ），我们指的是模 n 的整数环 \(\mathbf{Z}/n\mathbf{Z}\) 中可逆元素的个数。根据下一个命题， \(\varphi(n)\) 也是与 n 互素且满足 \(0 \leqslant k < n\) 的整数 k 的个数。

命题3. --- 设 k 和 \(n \geqslant 1\) 为两个整数。则以下断言等价：

\begin{enumerate}
\def\labelenumi{\alph{enumi})}
\item
  k 模 n 的剩余类在环 Z/nZ 中是可逆的；
\item
  k 模 n 的剩余类生成循环群 Z/nZ；
\item
  整数 \(k\) 和 \(n\) 互素（I，第112页）。
\end{enumerate}

条件 a) 和 b) 各自意味着存在一个整数 x，使得 \(kx \equiv 1 \pmod{n}\) ，即存在两个整数 x 和 y 使得 \(kx + ny = 1\) 。后一条件恰好意味着 k 与 n 互素。

推论1. --- 设 G 为 n 阶循环群，且 d 为 n 的一个因子。则 G 中阶为 d 的元素的个数等于 \(\varphi(d)\) 。特别地， \(\varphi(n)\) 是 G 的生成元的个数。

由于 G 同构于 Z/nZ，由命题3，G 的生成元的个数等于 \(\varphi(n)\) 。设 g 为 G 的一个生成元；则 G 中满足 \(h^{d}=1\) 的元素 h 构成 G 的由 \(g^{n/d}\) 生成的子群 H；该群是 d 阶循环群，且 G 中 d 阶元素恰为 H 的生成元，因此其个数为 \(\varphi(d)\) 。

推论2. --- 对每个整数 \(n \geqslant 1\) ，我们有

\[
\sum_ {d \mid n} \varphi (d) = n,\tag{1}
\]

其中整数 \(d\) 遍历所有 \(>0\) 的、 \(n\) 的因数组成的集合 \(^{1}\) 。

利用推论1的记号，G 的每个元素都有有限阶，该阶是 n 的一个因数 d，且对固定的 d，这样的元素有 \(\varphi(d)\) 个。

\(\varphi(n)\) 的计算基于以下两个公式：

\begin{enumerate}
\def\labelenumi{(\arabic{enumi})}
\setcounter{enumi}{1}
\item
  \(\varphi(mn) = \varphi(m)\varphi(n)\) （若 \(m\) 和 \(n\) 互素），
\item
  \(\varphi(p^a) = p^{a-1}(p - 1)\) （ \(p\) 素数， \(a \geqslant 1\) ）。
\end{enumerate}

第一个公式立即可由环 \(\mathbf{Z} / mn\mathbf{Z}\) 和 \((\mathbf{Z} / m\mathbf{Z}) \times (\mathbf{Z} / n\mathbf{Z})\) 同构（I，第112页）以及对两个环 A 和 B 有 \((\mathbf{A} \times \mathbf{B})^{*} = \mathbf{A}^{*} \times \mathbf{B}^{*}\) 这一事实得出。为证（3），我们注意到 \(p^{a}\) 的正因数是 1， \(p\) , \(p^{2}, \ldots, p^{a}\) ；因此整数 \(k\) 与 \(p^{a}\) 除 1 外无公因数，当且仅当它不被 \(p\) 整除；由于 \(p^{a-1}\) 个 \(p\) 的倍数位于 0 与 \(p^{a}-1\) 之间，我们确实得到（3）。

公式（2）和（3）立即表明

\[
\varphi (n) = n \prod_ {p} (1 - 1 / p),\tag{4}
\]

其中 p 遍历 n 的素因子集合。

一个 n 次单位根若其阶为 n，则称为本原的；若存在这样一个根 \(\zeta\) ，则群 \(\mu_{n}(\mathbf{K})\) 的阶为 n，且由 \(\zeta\) 生成。* 例如，C 中的 n 次本原单位根是数 \(e^{2\pi ik/n}\) ，其中 \(0 \leqslant k < n\) 且 k 与 n 互素。* 命题3的推论1现在蕴含以下结果。

命题4. --- 设 \(n \geqslant 1\) 为一个整数；假设存在 \(n\) 个 \(n\) 次单位根于 \(K\) （例如当 \(K\) 是可分封闭的且 \(n \cdot 1_{K} \neq 0\) 时即如此）。则 \(n\) 次本原单位根（ \(K\) 中的）的个数等于 \(\varphi(n)\) 。

